\title{CasePath: An Agentic, Process-First Architecture for Determining Evidence Requirements}
\author{Anonymous authors}
\maketitle

\begin{abstract}
Insurance claims are a useful testbed for AI decision systems because the evidence needed to resolve a case depends on where that case lies in a rule-governed process. A single fact can activate a different branch, create new obligations, and change which documents are justified. Yet many LLM-based approaches bypass this structure and predict document checklists directly from case text. Such requests may look reasonable while being weakly grounded: the system can ask for the right-looking document for the wrong reason. We introduce \casepath, an agentic, process-first architecture in which the source-defined process—not the language model—controls evidence acquisition. Agents read authoritative sources, place the case on the applicable process path, identify active or unresolved requirements, and request only the evidence needed to satisfy them. As evidence arrives, the case state is updated and the required checklist changes accordingly. We test this dependency using matched cases that differ in one branch-defining fact. Across 27 held-out pairs, \casepath makes 15 false-positive document changes, versus at least 27 for every comparison system. This turns document prediction into evidence planning that is selective, traceable, and auditable. \end{abstract}
